\thesisAppendix{В}{Скриншоты веб-составляющей}

В настоящем приложении приведены скриншоты основных пользовательских экранов веб-составляющей системы <<ВКурсе>>. Веб-составляющая объединяет три функциональных контура --- публичный лендинг с информацией о проекте, веб-приложение просмотра расписания и закрытый веб-сервис администрирования; рисунки сгруппированы соответственно этим контурам. В отличие от приложений~А и~Б, в которых композитные рисунки построены сеткой $2 \times 2$ из мобильных скриншотов, в настоящем приложении каждый рисунок представляет один экран целиком (за исключением рис.~\ref{fig:app-web-schedule-mobile}, где приведены два состояния мобильной раскладки).

\begin{figure}[h]
  \centering
  % TODO: заменить \fbox на \includegraphics{...} с реальным скриншотом главной страницы лендинга.
  \fbox{\parbox[c][6cm][c]{0.85\textwidth}{\centering\itshape Скриншот: главная страница лендинга \texttt{vcourse.app} --- ключевое сообщение, демонстрация мобильных приложений, призывы к действию}}
  \caption{Главная (приветственная) страница лендинга}
  \label{fig:app-web-landing}
\end{figure}

\begin{figure}[h]
  \centering
  % TODO: заменить \fbox на \includegraphics{...} с реальным скриншотом страницы <<О нас>>.
  \fbox{\parbox[c][6cm][c]{0.85\textwidth}{\centering\itshape Скриншот: страница <<О нас>> лендинга --- описание проекта, состав команды, контакты}}
  \caption{Страница <<О нас>> лендинга}
  \label{fig:app-web-about}
\end{figure}

\begin{figure}[h]
  \centering
  % TODO: заменить \fbox на \includegraphics{...} с реальным скриншотом страницы <<Блог>>.
  \fbox{\parbox[c][6cm][c]{0.85\textwidth}{\centering\itshape Скриншот: страница <<Блог>> лендинга --- лента записей с обложками, заголовками и кратким описанием}}
  \caption{Страница <<Блог>> лендинга}
  \label{fig:app-web-blog}
\end{figure}

\begin{figure}[h]
  \centering
  % TODO: заменить \fbox на \includegraphics{...} с реальным скриншотом страницы <<Университетам>>.
  \fbox{\parbox[c][6cm][c]{0.85\textwidth}{\centering\itshape Скриншот: страница <<Университетам>> лендинга --- предложение для учебных заведений, описание сервиса администрирования}}
  \caption{Страница <<Университетам>> лендинга}
  \label{fig:app-web-universities}
\end{figure}

\begin{figure}[h]
  \centering
  % TODO: заменить \fbox на \includegraphics{...} с реальным скриншотом десктопной раскладки.
  \fbox{\parbox[c][6cm][c]{0.85\textwidth}{\centering\itshape Скриншот: десктопная раскладка веб-приложения просмотра расписания --- постоянная боковая панель с выбором учебного заведения и точки зрения, центральная область с сеткой расписания}}
  \caption{Веб-приложение просмотра расписания: десктопная раскладка расписания учебной группы}
  \label{fig:app-web-schedule-desktop}
\end{figure}

\begin{figure}[h]
  \centering
  % TODO: заменить каждый \fbox на \includegraphics{...} с реальным скриншотом мобильной раскладки.
  \begin{minipage}[t]{0.40\textwidth}
    \centering
    \fbox{\parbox[c][9cm][c]{0.95\linewidth}{\centering\itshape Скриншот: мобильная раскладка веб-приложения просмотра расписания --- компактная сетка с свёрнутой боковой панелью}}\\[3pt]
    \small (а)
  \end{minipage}\hspace{0.05\textwidth}%
  \begin{minipage}[t]{0.40\textwidth}
    \centering
    \fbox{\parbox[c][9cm][c]{0.95\linewidth}{\centering\itshape Скриншот: мобильная раскладка веб-приложения просмотра расписания с выдвинутым меню \texttt{Sheet} (выбор справочной сущности)}}\\[3pt]
    \small (б)
  \end{minipage}
  \caption{Веб-приложение просмотра расписания: мобильная раскладка --- (а)~компактный вид со свёрнутой боковой панелью; (б)~вид с выдвинутым меню \texttt{Sheet} для выбора справочной сущности}
  \label{fig:app-web-schedule-mobile}
\end{figure}

\begin{figure}[h]
  \centering
  % TODO: заменить \fbox на \includegraphics{...} с реальным скриншотом диалога подробной информации о паре.
  \fbox{\parbox[c][6cm][c]{0.85\textwidth}{\centering\itshape Скриншот: диалог подробной информации о занятии, открываемый по нажатию на карточку пары в сетке расписания --- дисциплина, тип, участники, аудитория, заметки}}
  \caption{Веб-приложение просмотра расписания: диалог подробной информации о занятии}
  \label{fig:app-web-schedule-lesson-dialog}
\end{figure}

\begin{figure}[h]
  \centering
  % TODO: заменить \fbox на \includegraphics{...} с реальным скриншотом диалога поиска сущности.
  \fbox{\parbox[c][6cm][c]{0.85\textwidth}{\centering\itshape Скриншот: диалог поиска справочной сущности --- поле поискового запроса, переключатель типа сущности (группа, преподаватель, аудитория), список результатов}}
  \caption{Веб-приложение просмотра расписания: диалог поиска справочной сущности для открытия её расписания}
  \label{fig:app-web-schedule-search-dialog}
\end{figure}

\begin{figure}[h]
  \centering
  % TODO: заменить \fbox на \includegraphics{...} с реальным скриншотом экрана аутентификации.
  \fbox{\parbox[c][6cm][c]{0.85\textwidth}{\centering\itshape Скриншот: экран аутентификации сотрудника учебного заведения --- форма входа (логин и пароль), описание сервиса}}
  \caption{Веб-сервис администрирования: экран аутентификации}
  \label{fig:app-web-admin-auth}
\end{figure}

\begin{figure}[h]
  \centering
  % TODO: заменить \fbox на \includegraphics{...} с реальным скриншотом первого шага онбординга.
  \fbox{\parbox[c][6cm][c]{0.85\textwidth}{\centering\itshape Скриншот: первый шаг мастера онбординга --- форма заполнения информации об учебном заведении (название, адрес, логотип, реквизиты)}}
  \caption{Веб-сервис администрирования: первый шаг мастера онбординга --- заполнение информации об учебном заведении}
  \label{fig:app-web-admin-onboarding}
\end{figure}

\begin{figure}[h]
  \centering
  % TODO: заменить \fbox на \includegraphics{...} с реальным скриншотом главного дашборда.
  \fbox{\parbox[c][6cm][c]{0.85\textwidth}{\centering\itshape Скриншот: главный дашборд веб-сервиса администрирования --- трёхколоночный макет shadcn/ui, левая боковая панель разделов, центральная область со сводными показателями, правая контекстная панель}}
  \caption{Веб-сервис администрирования: главный дашборд сотрудника}
  \label{fig:app-web-admin-dashboard}
\end{figure}

\begin{figure}[h]
  \centering
  % TODO: заменить \fbox на \includegraphics{...} с реальным скриншотом справочника <<Группы>>.
  \fbox{\parbox[c][6cm][c]{0.85\textwidth}{\centering\itshape Скриншот: справочник учебных групп --- таблица TanStack Table с фильтрацией по факультету, форме обучения и году набора, действия редактирования и удаления}}
  \caption{Веб-сервис администрирования: справочник <<Группы>>}
  \label{fig:app-web-admin-groups}
\end{figure}

\begin{figure}[h]
  \centering
  % TODO: заменить \fbox на \includegraphics{...} с реальным скриншотом справочника <<Факультеты>>.
  \fbox{\parbox[c][6cm][c]{0.85\textwidth}{\centering\itshape Скриншот: справочник факультетов --- таблица TanStack Table с раскрываемыми строками для отображения кафедр в составе факультета}}
  \caption{Веб-сервис администрирования: справочник <<Факультеты>>}
  \label{fig:app-web-admin-faculties}
\end{figure}

\begin{figure}[h]
  \centering
  % TODO: заменить \fbox на \includegraphics{...} с реальным скриншотом диалога добавления факультета.
  \fbox{\parbox[c][6cm][c]{0.85\textwidth}{\centering\itshape Скриншот: модальный диалог добавления нового факультета --- поля названия, сокращённого названия, привязки к учебному заведению}}
  \caption{Веб-сервис администрирования: диалог добавления нового факультета на экране справочника <<Факультеты>>}
  \label{fig:app-web-admin-faculty-add}
\end{figure}

\begin{figure}[h]
  \centering
  % TODO: заменить \fbox на \includegraphics{...} с реальным скриншотом редактора расписания.
  \fbox{\parbox[c][6cm][c]{0.85\textwidth}{\centering\itshape Скриншот: редактор расписания --- сетка недели с карточками занятий учебной группы, индикаторы изменённых атомов, кнопки публикации и отката, всплывающие подсказки}}
  \caption{Веб-сервис администрирования: редактор расписания для учебной группы}
  \label{fig:app-web-admin-editor}
\end{figure}

\begin{figure}[h]
  \centering
  % TODO: заменить \fbox на \includegraphics{...} с реальным скриншотом диалога сохранения изменений.
  \fbox{\parbox[c][6cm][c]{0.85\textwidth}{\centering\itshape Скриншот: модальный диалог сохранения изменений в редакторе расписания --- сводный перечень накопленных дельт атомов, результаты серверной валидации ресурсных коллизий, кнопка публикации}}
  \caption{Веб-сервис администрирования: диалог сохранения изменений в редакторе расписания}
  \label{fig:app-web-admin-editor-save-dialog}
\end{figure}

\begin{figure}[h]
  \centering
  % TODO: заменить \fbox на \includegraphics{...} с реальным скриншотом редактора в тёмной теме и на английском.
  \fbox{\parbox[c][6cm][c]{0.85\textwidth}{\centering\itshape Скриншот: редактор расписания в тёмной теме интерфейса с английским языком --- иллюстрация поддержки темизации и локализации веб-сервиса администрирования}}
  \caption{Веб-сервис администрирования: редактор расписания в тёмной теме на английском языке}
  \label{fig:app-web-admin-editor-dark-en}
\end{figure}
